\section{Conclusion}

The pure-row presentation has an exact semantic decomposition. First take
the least carrier congruence for which every supplied row becomes
homomorphically coherent and kernel-stable. Then the ambient extension of
each row kernel determines precisely which remaining inputs acquire
carrier values. The resulting quotient and domains are consequences of
the equations, without any choice of a completion.

The local-to-global argument also determines the entire equality relation
on the named interface. Its horizon-one quotient records value coherence;
its horizon-two quotient records the full row geometry. Their class counts
give an exact stabilization test, while the sharp feedback family shows
why this proof-depth bound does not bound the number of structural updates.

Completion asks for an additional property: a retraction of the row
pushout onto the distinguished carrier. Its quotient solutions need not
have a least element and need not survive arbitrary coarsening. Their
intersection has a separate meaning as the carrier equality forced by
models with no external named action values. The strict repair-kernel gap
makes the difference from unrestricted initial semantics explicit.

Additional composition and group laws connect rows and require their own
full-theory quotient. One-sorted coupling can introduce mixed terms or
reconstruct operations, depending on the equations imposed. These
boundaries specify where the pure-row characterization applies and which
new interactions a more general algorithm must account for. They also
provide a common foundation for the companion studies of equality
visibility and certificate-producing computation.
